\documentclass[11pt,a4paper]{article}
\usepackage[margin=2.5cm]{geometry}
\usepackage{graphicx}
\usepackage{booktabs}
\usepackage{longtable}
\usepackage{xcolor}
\usepackage{hyperref}
\usepackage{amsmath}
\usepackage{natbib}
\usepackage{caption}
\usepackage{array}
\usepackage{multirow}
\usepackage{float}

\hypersetup{colorlinks=true, linkcolor=blue!60!black, citecolor=blue!60!black, urlcolor=blue!60!black}

\title{Binding QTL Act as Quantitative Modulators of Gene Regulation\\Through Motif Redundancy and Combinatorial TF Disruption in Maize}
\author{Michael Banf$^{1}$, Thomas Hartwig$^{2,*}$\\[6pt]
\small $^1$Perelyn, Munich, Germany\\
\small $^2$Institute of Plant Genetics, Heinrich Heine University D\"usseldorf / MPI for Plant Breeding Research / CEPLAS\\
\small $^*$Corresponding author: hartwit@hhu.de}
\date{}

\begin{document}
\maketitle

% ============================================================
\begin{abstract}
Binding quantitative trait loci (bQTL) identify cis-regulatory variants that alter transcription factor occupancy, but whether they act as binary switches or quantitative modulators of gene expression remains unexplored. Here we systematically characterize the motif landscape of 147,942 bQTL from 25 maize F1 hybrids, linking motif redundancy to allele-specific expression (ASE) across 13,895 genes. We find that 95.8\% of bQTL-disrupted TF binding motifs are redundant---the same TF family retains additional binding sites within $\pm$200~bp---establishing that bQTL are quantitative modulators rather than binary switches. The buffering effect is TF-family-specific: MYB ($\rho = -0.049$, $p = 6.7 \times 10^{-11}$), MYB\_related ($\rho = -0.053$, $p = 1.9 \times 10^{-11}$), and Dof ($\rho = -0.048$, $p = 2.2 \times 10^{-8}$) motifs show strong buffering, while SBP ($\rho = +0.063$, $p = 6.9 \times 10^{-9}$), bHLH, bZIP, NAC, and ARF motifs are sensitive to disruption. This defines a two-tier regulatory architecture: abundant, redundant ``buffered'' families that tolerate variant disruption versus sparse ``sensitive'' families where single-variant loss has disproportionate effects. Combinatorial disruption across multiple TF families is the strongest predictor of allelic expression magnitude ($\rho = +0.063$, $p = 1.2 \times 10^{-13}$), with genes disrupted in all 18 families showing 24\% higher allelic imbalance than single-family disrupted genes. Integration with single-cell chromatin and transcriptomic data resolves 271 bQTL to three leaf cell types, revealing cell-type-specific regulatory chains including a novel HSFA6a--NAC73 feedback loop in bulliform cells that couples drought sensing to cell wall remodeling. These findings redefine bQTL as quantitative fine-tuning variants embedded in a robust, buffered regulatory architecture, with implications for understanding cis-regulatory evolution and precision crop improvement.
\end{abstract}

% ============================================================
\section{Introduction}

Natural genetic variation in transcription factor (TF) binding is a primary driver of phenotypic diversity. In maize, the leaf pan-cistrome---generated by MNase-defined Occupancy Analysis of sequencing (MOA-seq) across 25 F1 hybrids of B73 with NAM founder lines---identified 147,942 binding quantitative trait loci (bQTL) that collectively explain $>$72\% of heritable trait variation across 143 agronomic phenotypes \citep{engelhorn2025}. Each bQTL marks a position where TF occupancy differs between the B73 and NAM alleles within the same nucleus, establishing cis-acting regulatory variation.

The original characterization of these bQTL treated them as a population-level statistical phenomenon: positions were identified, classified by mechanism (genotype-dependent, methylation-dependent, or interaction), and linked to distant phenotypes through mediation analysis. However, several fundamental mechanistic questions remained unaddressed:

\begin{enumerate}
    \item \textbf{Which TF} is affected at each bQTL? MOA-seq detects protein occupancy changes without identifying the bound factor.
    \item Does the bQTL disrupt the \textbf{only} binding site for that TF, or one of many?
    \item Do bQTL act as \textbf{binary switches} (on/off) or \textbf{quantitative dials} (fine-tuning)?
    \item What is the \textbf{motif landscape} surrounding each bQTL?
    \item How does \textbf{motif redundancy} buffer against single-variant effects?
    \item In which \textbf{cell types} do specific bQTL exert their regulatory effects?
\end{enumerate}

These questions are critical for understanding the mechanistic basis of cis-regulatory variation and its translation to phenotype. The F1 hybrid system provides an ideal natural experiment: both alleles reside in the same nucleus, exposed to identical trans factors. Any difference in TF occupancy or gene expression must be cis-encoded in the DNA sequence. Allele-specific expression (ASE) measurements from these hybrids directly quantify the functional consequence of each regulatory variant---and because both alleles typically encode identical proteins, the ASE signal reflects a pure \textit{dosage} change rather than a qualitative functional difference.

Large-scale synthetic promoter studies have established that transcription is quantitatively modulated by motif multiplicity. In yeast, expression levels across 100 million random promoters scale directly with the number and identity of embedded TF binding motifs \citep{deboer2020}. In plants, massively parallel reporter assays (MPRA) confirm that promoter strength correlates with motif composition in maize, Arabidopsis, and sorghum \citep{jores2021}. Homotypic clusters of TF binding sites (HCTs)---multiple copies of the same motif in a promoter---are a key architectural feature of eukaryotic enhancers and promoters, providing cooperative binding, temporal dynamics, and robustness against single-motif perturbation \citep{gotea2010}.

These precedents predict that most natural cis-regulatory variants should act as quantitative modulators rather than binary switches. We test this hypothesis systematically across the maize bQTL landscape by (i) scanning all bQTL-containing promoters for TF motif disruption and redundancy, (ii) linking motif context to measured allelic expression differences, (iii) identifying family-specific buffering and sensitivity patterns, and (iv) resolving regulatory chains to individual cell types through integration with single-cell data.

% ============================================================
\section{Results}

\subsection{Most bQTL disrupt redundant TF binding motifs}

We mapped all 147,942 bQTL to gene promoters (2~kb upstream of TSS) using the B73 NAM5 reference annotation, identifying 33,634 unique bQTL positions in the promoters of 13,895 genes with allele-specific expression data. For each promoter bQTL, we performed a variant-overlap TF motif scan: a $\pm$28~bp window around the variant was extracted and scored against 222 position weight matrices (PWMs) from PlantTFDB 5.0, representing 18 TF families (log-odds threshold $\geq 6.0$). Only motifs whose match interval physically overlaps the variant nucleotide were retained, ensuring that each hit represents a motif directly disrupted by the variant.

This scan produced 182,401 bQTL-gene-TF triples (unique combinations of variant position, target gene, and TF family with a disrupted motif). On average, each bQTL disrupts motifs from 5.4 TF families, consistent with the high density of overlapping binding sites in plant promoters.

For each disrupted motif, we quantified \textbf{motif redundancy}: the number of additional intact binding sites for the same TF family within $\pm$200~bp of the variant. Strikingly, \textbf{95.8\% of disrupted motifs are redundant}---the promoter retains at least one additional binding site for the same TF family. Only 4.2\% (7,621/182,401) of disrupted motifs are ``sole'' motifs, representing the only binding opportunity for that TF family in the local regulatory window (Figure~\ref{fig:redundancy}A).

This establishes that the vast majority of bQTL reduce, but do not eliminate, TF binding capacity at their target promoters. Rather than acting as binary on/off switches, bQTL function as quantitative modulators that fine-tune regulatory input within a buffered motif architecture.

\subsection{Motif redundancy buffers allelic expression differences: a TF-family-specific effect}

If motif redundancy provides functional buffering, genes where a bQTL disrupts one of many same-family motifs should show smaller allelic expression differences than genes where the bQTL disrupts the sole motif. We tested this by correlating motif redundancy with the absolute allelic expression ratio $|\log_2(\text{B73/NAM})|$ from the Engelhorn et al.\ ASE data (23,807 genes, 24 hybrids).

The overall correlation is weak and marginally significant ($\rho = -0.004$, $p = 0.089$), reflecting the fact that TF families differ fundamentally in their regulatory logic. When analyzed per family, a striking pattern emerges (Table~\ref{tab:families}, Figure~\ref{fig:redundancy}C):

\textbf{Buffered families} (negative $\rho$: more redundancy $\rightarrow$ smaller ASE):
\begin{itemize}
    \item \textbf{MYB\_related}: $\rho = -0.053$, $p = 1.9 \times 10^{-11}$ (n = 15,872)
    \item \textbf{MYB}: $\rho = -0.049$, $p = 6.7 \times 10^{-11}$ (n = 17,803)
    \item \textbf{Dof}: $\rho = -0.048$, $p = 2.2 \times 10^{-8}$ (n = 13,566)
    \item \textbf{HD-ZIP}: $\rho = -0.022$, $p = 0.044$ (n = 8,672)
\end{itemize}

\textbf{Sensitive families} (positive $\rho$: more redundancy $\rightarrow$ \textit{larger} ASE):
\begin{itemize}
    \item \textbf{SBP}: $\rho = +0.063$, $p = 6.9 \times 10^{-9}$ (n = 8,501)
    \item \textbf{bHLH}: $\rho = +0.043$, $p = 1.5 \times 10^{-4}$ (n = 7,925)
    \item \textbf{ARF}: $\rho = +0.036$, $p = 6.5 \times 10^{-4}$ (n = 9,154)
    \item \textbf{NAC}: $\rho = +0.033$, $p = 1.2 \times 10^{-4}$ (n = 13,454)
    \item \textbf{bZIP}: $\rho = +0.034$, $p = 3.1 \times 10^{-3}$ (n = 7,832)
    \item \textbf{ERF}: $\rho = +0.024$, $p = 4.3 \times 10^{-3}$ (n = 14,465)
\end{itemize}

The MYB superfamily (MYB + MYB\_related, the most abundant TF families in maize promoters) provides the strongest evidence for quantitative buffering. MYB TFs are established regulators of cell wall biosynthesis, flavonoid metabolism, and stress responses. Their binding sites are characteristically abundant in plant promoters---our data shows a mean of 10--15 MYB family motifs per promoter---and the strong negative correlation confirms that this redundancy provides genuine functional buffering against individual variant disruption.

Conversely, the sensitive families (SBP, bHLH, bZIP, NAC, ARF) show the paradoxical pattern where more same-family motifs correlate with \textit{larger} allelic expression differences. This likely reflects combinatorial regulatory logic: these families often function in obligate heteromeric complexes (bHLH-bZIP, ARF-auxin response) or as dosage-sensitive developmental regulators (SBP, NAC), where motif multiplicity may indicate more complex regulation rather than simple buffering.

\subsection{Disruption fraction confirms the buffering model}

To quantify the proportion of regulatory capacity affected by each bQTL, we computed the \textbf{disruption fraction}: $1 / (n_\text{total motifs})$, representing the fraction of same-family binding sites disrupted. For buffered families, a higher disruption fraction should predict larger ASE:

\begin{itemize}
    \item \textbf{MYB\_related}: $\rho = +0.053$ ($p = 1.9 \times 10^{-11}$) --- confirmed
    \item \textbf{MYB}: $\rho = +0.049$ ($p = 6.7 \times 10^{-11}$) --- confirmed
    \item \textbf{Dof}: $\rho = +0.048$ ($p = 2.2 \times 10^{-8}$) --- confirmed
\end{itemize}

These are the mathematical complements of the redundancy correlations, but they provide the intuitive interpretation: when a bQTL disrupts a larger \textit{proportion} of MYB binding capacity (e.g., 1 of 2 sites vs.\ 1 of 20), the allelic expression difference is proportionally larger. This is the hallmark of a quantitative, dose-dependent regulatory system.

For the sensitive families, the disruption fraction shows negative correlations (SBP: $\rho = -0.063$, bHLH: $\rho = -0.043$), consistent with their non-additive regulatory logic.

\subsection{Combinatorial disruption across TF families is the strongest predictor of allelic expression}

Individual motif redundancy captures within-family buffering, but gene regulation involves multiple TF families acting in concert. We tested whether the \textbf{number of different TF families} with motifs disrupted at a gene's bQTL predicts the magnitude of its allelic expression difference.

This is the strongest effect in our analysis: the number of disrupted TF families correlates with $|\log_2(\text{B73/NAM})|$ at $\rho = +0.063$ ($p = 1.2 \times 10^{-13}$, n = 13,895 genes; Figure~\ref{fig:combinatorial}A). Genes with motifs from all 18 families disrupted show a median ASE of 0.633, compared with 0.509 for genes with motifs from only 1 family---a 24\% increase in allelic imbalance.

This combinatorial effect exceeds the individual family effects, suggesting that multi-family disruption creates regulatory perturbation that exceeds the sum of single-family effects. When a promoter variant simultaneously disrupts binding sites for multiple TF families, it affects diverse regulatory inputs (activation, repression, stress response, developmental patterning), collectively amplifying the allelic expression difference.

We also tested which specific families predict higher ASE when disrupted at a gene's promoter (Table~\ref{tab:family_disruption}). All 11 significant families show the same direction: disruption predicts \textit{higher} ASE. The largest effects are for ERF (median difference +0.074, $p = 6.0 \times 10^{-22}$), LBD (+0.056, $p = 8.2 \times 10^{-14}$), bZIP (+0.056, $p = 1.7 \times 10^{-11}$), and bHLH (+0.052, $p = 1.5 \times 10^{-11}$).

Notably, MYB is \textit{not} among the significant families for the disruption-vs-no-disruption comparison, despite being the most frequently disrupted family. This is consistent with its buffered status: because MYB motifs are so abundant and redundant, their disruption alone does not significantly alter gene expression. It is the disruption of sensitive families (ERF, bHLH, bZIP, SBP) that drives the allelic signal.

\subsection{Cross-hybrid consistency increases with combinatorial disruption}

The maize F1 hybrid panel provides ASE measurements across 24 different genetic backgrounds (B73 $\times$ 24 NAM founders). We tested whether combinatorial disruption predicts not only the magnitude but also the \textit{variability} of ASE across hybrids.

The number of disrupted TF families correlates with ASE standard deviation across hybrids ($\rho = +0.068$, $p = 2.2 \times 10^{-14}$, n = 12,644 genes). Genes where many TF families are disrupted show more variable expression across different hybrid backgrounds. This suggests that multi-family disruption creates genotype-dependent regulatory instability: the expression consequence of a promoter variant depends on which trans factors are available, which differs across NAM founder genotypes.

\subsection{TF expression level does not modulate motif disruption effects under baseline conditions}

The cross-hybrid variation in ASE raises a natural question: does TF abundance modulate the functional consequence of motif disruption? Under a TF saturation model, high TF concentrations would buffer against motif loss (excess TF molecules rebind remaining sites), while low TF concentrations would amplify the effect (each binding site is critical). We tested this directly by treating the 19 NAM hybrids as a natural TF dose series.

For each TF family, we computed per-hybrid expression from allele-specific read counts (median across family members) and fitted an interaction model: $|\text{ASE}| \sim \beta_1 \cdot \text{disrupted} + \beta_2 \cdot \log_2(\text{TF\_expr}) + \beta_3 \cdot \text{disrupted} \times \log_2(\text{TF\_expr})$, where a positive $\beta_3$ indicates sub-saturation (motif loss hurts more when TF is scarce) and a negative $\beta_3$ indicates buffering by TF abundance. The analysis covered 730,497 observations across 27,118 bQTL and 12,061 genes.

No TF family showed a significant interaction (all FDR $> 0.75$; global Wilcoxon $p = 0.99$; mean $\beta_3 = -0.003$). The effect of motif disruption on allelic expression is independent of TF expression level across the natural range of variation within the hybrid panel.

This result has two implications. First, it confirms that motif disruption effects are structurally encoded---determined by the \emph{cis}-regulatory architecture (redundancy, combinatorial context) rather than modulated by the \emph{trans} TF concentration landscape. This is consistent with TF families operating in the responsive regime ($C_\text{free} \approx K_d$) under well-watered conditions, where occupancy is sensitive to motif presence but approximately linear with TF concentration. Second, it motivates the hypothesis that condition-dependent bQTL switching requires large changes in TF levels---such as the 2--30$\times$ inductions observed under drought stress---rather than the modest variation ($\text{CV} \approx 0.58$) across hybrids within a single condition.

\subsection{A two-tier regulatory architecture}

The per-family results define a \textbf{two-tier regulatory architecture} in maize promoters (Figure~\ref{fig:combinatorial}B):

\textbf{Tier 1 --- Buffered families} (MYB, MYB\_related, Dof, HD-ZIP): Abundant binding sites provide redundancy. Disruption of a single motif has minimal effect on expression because multiple alternative sites remain. These families act as quantitative modulators---their motif count tunes expression level, but individual variants have small marginal effects. This is consistent with homotypic clustering as a buffering mechanism \citep{gotea2010}.

\textbf{Tier 2 --- Sensitive families} (SBP, bHLH, bZIP, NAC, ARF, ERF): Fewer motifs per promoter (higher fraction of sole motifs). Disruption has disproportionate effects on expression. These families often function in combinatorial complexes or as dosage-sensitive developmental switches, where motif multiplicity reflects regulatory complexity rather than redundancy.

This architecture has implications for cis-regulatory evolution: buffered family binding sites are free to accumulate variants with minimal phenotypic consequence (genetic drift), while sensitive family binding sites are under stronger purifying selection. Indeed, our earlier analysis of bQTL enrichment patterns showed that developmental TF families (HD-ZIP, WOX, Dof) have binding sites \textit{depleted} at bQTL genome-wide---consistent with purifying selection preserving these regulatory elements \citep{engelhorn2025}.

\subsection{Genome-wide motif count correlates with expression level}

To establish the quantitative relationship between motif multiplicity and gene expression independently of bQTL, we scanned 500~bp promoters of 5,000 stratified genes for TF motifs and correlated count with expression level (mean CPM).

Seven of twelve tested families showed significant correlations (Spearman, $p < 0.05$): five positive (HD-ZIP: $\rho = +0.049$, TCP: $+0.042$, G2-like: $+0.040$, MYB: $+0.039$, Dof: $+0.035$) and two negative (NAC: $-0.033$, SBP: $-0.028$). The positive correlations are consistent with additive transcriptional activation: more binding sites recruit more TF molecules, producing higher transcription. The negative correlations for NAC and SBP may reflect repressor functions or negative autoregulation.

The correlations are weak but significant, as expected: motif presence is necessary but not sufficient for TF binding (chromatin accessibility, cofactor availability, and post-translational modification all modulate actual occupancy). The key finding is directional: for the majority of TF families, more motifs correspond to more expression, supporting a quantitative model.

\subsection{Case studies: the motif landscape of functionally characterized bQTL}

\subsubsection{ZmTRE1: a heavily buffered MYB target}

ZmTRE1 (Zm00001eb359620) encodes trehalase, an enzyme in trehalose-6-phosphate signaling that regulates ABA-mediated drought responses. The Engelhorn et al.\ paper identified a Dof TF (ZmDOF17) binding this promoter via ChIP-seq. Our motif scan reveals a single bQTL (chr8:152367402, MetOnly mechanism) that disrupts motifs from 4 TF families: MYB (+21 additional sites), SBP (+25), ERF (+5), and ARF (+5). The exceptionally high MYB redundancy (21 additional sites) predicts strong buffering---and indeed, ZmTRE1 shows one of the lowest ASE values in our dataset ($|\log_2| = 0.285$), despite being a highly expressed gene (557 CPM). The variant fine-tunes rather than switches.

\subsubsection{HSFA6a: combinatorial disruption at a stress master regulator}

HSFA6a (Zm00001eb160080) encodes a heat stress transcription factor that shows 30-fold drought induction. Four bQTL in its promoter span 154~bp and collectively disrupt motifs from 14 TF families, including NAC, bHLH, bZIP, MYB, Dof, ERF, G2-like, TCP, SBP, MYB\_related, C2H2, MIKC\_MADS, GATA, and ARF. The bQTL types include MetP (methylation $\times$ phenotype), GenoP (genotype $\times$ phenotype), and MetOnly, indicating that both genetic and epigenetic variation modulate this promoter.

Despite this extensive combinatorial disruption, all motifs are redundant (no sole motifs)---the promoter is saturated with TF binding sites. The measured ASE ($|\log_2| = 0.676$) is above the genome-wide median (0.548), consistent with the combinatorial disruption model: many families disrupted, but each individually buffered.

\subsubsection{Polyol transporter 5: the most complex regulatory landscape}

Polyol transporter 5 (Zm00001eb166210) has 13 bQTL across its 2~kb promoter, disrupting motifs from 17 of 18 TF families. This gene shows the highest ASE in our case study set ($|\log_2| = 1.327$), consistent with the combinatorial disruption prediction. Three motifs are sole (bHLH, bZIP, TCP), while the majority are highly redundant (LBD: up to +35, WRKY: +27, MYB: +23, NAC: +23). The gene encodes an osmolyte transporter involved in drought osmotic adjustment in bulliform cells.

\subsubsection{GBF3: mixed buffering with sole sensitive motifs}

GBF3 (Zm00001eb157860) is a bZIP transcription factor that confers drought tolerance when overexpressed \citep{ramegowda2017}. Its promoter contains 5 bQTL disrupting 13 TF families. Critically, two motifs are sole: HSF and LBD. The sole HSF motif is particularly notable given the HSF paradox: HSF binding sites are depleted at bQTL genome-wide (purifying selection), yet here a sole HSF site exists at the GBF3 promoter variant. The measured ASE ($|\log_2| = 0.880$) is elevated, consistent with the presence of sole sensitive-family motifs.

% ============================================================
\subsection{Cell-type resolution of regulatory chains}

\subsubsection{Framework}

While the quantitative modulation analysis characterizes the bulk regulatory logic of all 147,942 bQTL, bulk tissue profiling cannot resolve which cell types are affected by each variant. To address this, we integrated the bQTL data with single-cell chromatin accessibility (scATAC-seq) and single-nucleus RNA-seq (snRNA-seq) from maize leaf \citep{marand2021}, constructing \textbf{complete regulatory chains}: bQTL $\rightarrow$ TF motif disrupted $\rightarrow$ TF expressed $\rightarrow$ target gene expressed, resolved per cell type.

Three cell types---bulliform, guard cell, and epidermis---have both scATAC and snRNA data, enabling complete chain reconstruction. Of 20,000 bQTL tested, 271 pass all evidence layers (chromatin accessibility + motif + TF expression + target expression) across the three cell types, with 103 forming fully complete chains.

\subsubsection{Bulliform cells dominate the regulatory landscape}

Cell-type assignment revealed a strong bias: 161 bQTL (59\%) assigned to bulliform cells, 75 (28\%) to guard cells, and 35 (13\%) to epidermis. Bulliform cells are specialized motor cells that drive leaf rolling---the primary drought avoidance mechanism in grasses---consistent with disproportionate regulatory variation in this stress-responsive cell type.

\subsubsection{HSFA6a--NAC73 feedback loop}

The most striking discovery from the cell-type analysis is a mutual regulatory loop between HSFA6a and NAC73 in bulliform cells. NAC motifs are disrupted at the HSFA6a bQTL, while HSF motifs are disrupted at the NAC73 bQTL. Both TFs are expressed in bulliform cells, and both target genes have complete chains.

This feedback couples ABA-mediated drought sensing (HSFA6a) with secondary cell wall remodeling (NAC73/SND2), the program that controls mechanical stiffness in bulliform cells. Switch-like behavior is predicted: once either TF exceeds a threshold under drought stress, reciprocal reinforcement rapidly commits bulliform cells to simultaneous stress activation and cell wall restructuring. The rice ortholog of NAC73 (OsSND2) directly controls leaf rolling through cell wall thickening \citep{ye2018}, providing cross-species validation of this mechanistic link.

A three-level cascade extends the network: PTM (a chloroplast envelope-bound Trihelix TF) $\rightarrow$ NAC73 $\rightarrow$ HSFA6a, tracing a complete arc from chloroplast stress perception through cell wall remodeling to the heat-stress program \citep{sun2011}.

\subsubsection{Guard cell targets for stomatal engineering}

Twenty-five bQTL have complete chains specifically in guard cells, identifying targets for stomatal engineering including an uncharacterized gene with 3-fold drought induction (Zm00001eb282870), a membrane Mg$^{2+}$ transporter (stomatal ion dynamics), and FATTY ACID EXPORT 3 (cuticular lipid export with WRKY regulation).

\subsubsection{MetOnly bQTL are enriched for drought-responsive targets}

MetOnly (methylation-dependent) bQTL show a higher rate of complete chain formation (49.4\%) than GenoOnly (34.8\%) and are enriched for drought-upregulated targets (26\% vs.\ 11\%). This connects epigenetic variation disproportionately to stress-responsive programs, consistent with DNA methylation as a stress-responsive epigenetic mark in plants.

% ============================================================
\section{Discussion}

\subsection{bQTL as fine-tuning variants in a buffered regulatory architecture}

Our systematic analysis of 33,634 promoter bQTL establishes that the vast majority of natural cis-regulatory variants in maize act as quantitative modulators of gene expression rather than binary switches. The 95.8\% motif redundancy rate indicates that maize promoters are architecturally buffered against single-variant disruption of TF binding---a design principle consistent with the evolutionary robustness of cis-regulatory elements.

This buffering is not uniform. The two-tier architecture we describe---where MYB/Dof/HD-ZIP sites provide redundant quantitative buffering while SBP/bHLH/bZIP/NAC/ARF sites are sensitive to disruption---creates a regulatory system where some regulatory inputs (general activation, cell-type maintenance) are robust to sequence variation, while others (developmental decisions, stress responses) are sensitive. This differential sensitivity may explain why bQTL can explain 72\% of heritable trait variation despite individually having small effects: the sensitive-family disruptions provide the variance, while the buffered families maintain baseline regulatory programs.

\subsection{Combinatorial disruption as the dominant regulatory mechanism}

The finding that combinatorial disruption (number of TF families affected) is the strongest predictor of allelic expression ($\rho = +0.063$, $p = 1.2 \times 10^{-13}$) has important implications. Single-motif disruption analyses---the standard approach in regulatory genomics---may underestimate the functional impact of cis-regulatory variants. The relevant unit of regulatory perturbation is not the individual binding site but the multi-family regulatory context of the promoter.

This is consistent with the combinatorial TF binding model \citep{spitz2012}, where gene expression is determined by the integrated input of multiple TF families acting through distinct binding sites. A promoter variant that simultaneously disrupts binding for activators, repressors, and stress-responsive factors creates a regulatory perturbation qualitatively different from single-family disruption.

\subsection{The quantitative modulation model and synthetic biology}

Our findings directly inform synthetic promoter design. The de Boer et al.\ \citep{deboer2020} model of expression as a function of motif count is validated in a natural context: maize promoters follow the same quantitative logic as synthetic yeast promoters. This suggests that rational promoter engineering---adding or removing specific TF binding sites---should produce predictable changes in expression level, with MYB/Dof sites providing fine-tuning resolution and SBP/bHLH sites providing larger expression shifts per variant.

The cell-type-resolved regulatory chains further enable precision engineering: modifying specific TF motifs at bQTL positions in bulliform cell promoters (e.g., NAC73, polyol transporters) should affect leaf rolling and osmotic regulation without pleiotropic effects on other cell types.

\subsection{Limitations}

Several limitations should be noted. First, PWM-based motif scanning detects potential binding sites based on sequence match, not actual TF occupancy. Not all scanned motifs are functionally bound under physiological conditions. Second, the ASE data represents a population average across multiple cell types, while individual bQTL may act in specific cell types. Third, the $\pm$200~bp redundancy window is arbitrary; true functional redundancy may extend further or be constrained by nucleosome positioning. Fourth, the direction of effect (whether a variant increases or decreases binding) cannot be determined from the bQTL data, which reports only the presence of allelic occupancy differences. Finally, the cross-stage mapping between V6 bQTL and V2 single-cell data, while justified for stable regulatory elements, may miss stage-specific effects.

% ============================================================
\section{Methods}

\subsection{Data sources}

\textbf{bQTL.} 147,942 binding QTL from 25 maize F1 hybrids (B73 $\times$ NAM founders), from Engelhorn et al.\ \citep{engelhorn2025} Supplementary Table S6. Classified by mechanism: GenoOnly, MetOnly, GenoP, MetP.

\textbf{Allele-specific expression.} Haplotype-specific read counts from Engelhorn et al.\ Table S13a, providing ASE for 23,807 genes across 24 hybrids (well-watered condition). Gene-level summary: mean $|\log_2(\text{B73/NAM})|$ across hybrids.

\textbf{Expression and drought response.} Bulk RNA-seq expression (CPM) and drought fold-change from Engelhorn et al.\ for 24,796 genes.

\textbf{TF motifs.} 222 position weight matrices from PlantTFDB 5.0, representing 18 TF families in \textit{Zea mays}.

\textbf{Gene annotation.} B73 NAM5 reference genome (GCA\_902167145.1) and GFF3 annotation from MaizeGDB.

\textbf{Single-cell data.} scATAC-seq and snRNA-seq from maize leaf \citep{marand2021}, covering six and four cell types, respectively.

\subsection{Motif redundancy analysis}

For each bQTL in a gene promoter (2~kb upstream of TSS), a variant-overlap scan extracted $\pm$28~bp (maximum PWM width) around the variant and scored against all 222 PWMs using log-odds scoring with maize-specific background frequencies ($f_A = f_T = 0.264$, $f_C = f_G = 0.236$). Both forward and reverse complement strands were scored; hits above a log-odds threshold of $\geq$6.0 that physically overlap the variant nucleotide were classified as ``disrupted motifs.''

For each disrupted motif, redundancy was assessed by scanning $\pm$200~bp flanking the variant for additional intact binding sites of the same TF family (same log-odds threshold). The number of additional hits constitutes the redundancy count; motifs with zero additional hits are classified as ``sole.''

\subsection{Disruption fraction}

The disruption fraction for each bQTL-gene-TF triple was computed as $1 / (n_\text{other} + 1)$, where $n_\text{other}$ is the redundancy count. This represents the proportion of same-family binding capacity disrupted by the variant.

\subsection{Combinatorial disruption}

For each gene, the number of distinct TF families with disrupted motifs across all its promoter bQTL was computed. This gene-level metric captures the breadth of regulatory perturbation.

\subsection{Statistical analysis}

All correlations use Spearman rank correlation (non-parametric). Group comparisons use Mann-Whitney U tests. Per-family analyses include all families with $n \geq 50$ triples. Multiple testing correction is noted where applicable; primary results (per-family correlations) survive Bonferroni correction across 18 families.

\subsection{TF dose--disruption interaction model}

To test whether TF expression level modulates the effect of motif disruption, we constructed a per-hybrid analysis table. NAM founder genotypes at all 147,942 bQTL positions were obtained by mapping 201~bp B73 windows ($\pm$100~bp around each variant) to each of 19 NAM founder assemblies (Hufford et al.\ 2021) using \texttt{minimap2 -cx sr --cs} and parsing the cs alignment tags (union coverage: 99.6\% of bQTL). Per-hybrid TF family expression was computed as the median total read count (B73~+~NAM allele) across all expressed genes of each family (1,413 TF genes from 27 families, identified via EntAP functional annotations). For each combination of bQTL, target gene, TF family, and hybrid, we recorded: (i)~whether the founder allele disrupts the TF motif at the variant (from genotype), (ii)~the absolute ASE effect ($|\log_2(\text{B73/NAM})|$, minimum 10 total reads), and (iii)~the TF family expression in that hybrid. The interaction model $|\text{ASE}| = \beta_0 + \beta_1 \cdot \text{disrupted} + \beta_2 \cdot z(\log_2 \text{TF\_expr}) + \beta_3 \cdot \text{disrupted} \times z(\log_2 \text{TF\_expr}) + \varepsilon$ was fitted per family by OLS, with TF expression $z$-scored within each family for comparability. FDR correction (Benjamini--Hochberg) was applied across families.

\subsection{Cell-type regulatory chain construction}

Each bQTL position was intersected with cell-type-specific accessible chromatin regions (ACRs) from scATAC-seq. For bQTL within ACRs, TF motif scanning identified disrupted motifs ($\pm$28~bp, threshold $\geq$6.0). Complete chains (score = 4) required: ACR overlap + motif present + TF family expressed $>$10 CPM in that cell type + target gene expressed $>$10 CPM.

\subsection{Code and data availability}

All analysis scripts (44\_full\_quantitative\_modulation.py, 45\_case\_studies.py), raw results (quantitative\_modulation\_full.csv, per\_family\_redundancy.csv, gene\_level\_modulation.csv), and figures are available at [repository URL].

% ============================================================
\section{Conclusions}

We establish that bQTL in maize act primarily as quantitative modulators of gene expression rather than binary switches: 95.8\% of disrupted TF binding motifs are redundant, and the allelic expression consequence is determined by the proportion of same-family binding capacity affected. This quantitative modulation is TF-family-specific, defining a two-tier regulatory architecture where MYB/Dof sites are buffered and SBP/bHLH/bZIP/NAC sites are sensitive to disruption. Combinatorial disruption across multiple TF families is the strongest predictor of allelic expression magnitude ($\rho = +0.063$, $p = 1.2 \times 10^{-13}$), establishing multi-family perturbation as the dominant mechanism of cis-regulatory variation. Cell-type-resolved regulatory chains reveal that this quantitative variation has cell-type-specific consequences, including a novel HSFA6a--NAC73 feedback loop coupling drought sensing to cell wall remodeling in bulliform cells. These findings provide a mechanistic framework for understanding how natural cis-regulatory variation produces the quantitative phenotypic diversity that underlies crop adaptation and improvement.

% ============================================================

\begin{table}[!ht]
\centering
\caption{\textbf{Per-family motif redundancy vs.\ allelic expression.} Spearman correlation between the number of additional same-family motifs ($\pm$200~bp) and $|\log_2(\text{B73/NAM})|$ for each TF family. Direction: Buffered = negative $\rho$ (more motifs $\rightarrow$ smaller ASE); Sensitive = positive $\rho$.}
\label{tab:families}
\small
\begin{tabular}{lrrrrll}
\toprule
\textbf{Family} & \textbf{n} & \textbf{\% sole} & \textbf{$\rho$} & \textbf{$p$} & \textbf{Sig.} & \textbf{Direction} \\
\midrule
MYB\_related & 15,872 & 1.0\% & $-$0.053 & $1.9 \times 10^{-11}$ & *** & Buffered \\
MYB & 17,803 & 0.3\% & $-$0.049 & $6.7 \times 10^{-11}$ & *** & Buffered \\
Dof & 13,566 & 1.6\% & $-$0.048 & $2.2 \times 10^{-8}$ & *** & Buffered \\
HD-ZIP & 8,672 & 6.4\% & $-$0.022 & $4.4 \times 10^{-2}$ & * & Buffered \\
\midrule
ERF & 14,465 & 2.2\% & $+$0.024 & $4.3 \times 10^{-3}$ & ** & Sensitive \\
NAC & 13,454 & 1.7\% & $+$0.033 & $1.2 \times 10^{-4}$ & *** & Sensitive \\
bZIP & 7,832 & 11.5\% & $+$0.034 & $3.1 \times 10^{-3}$ & ** & Sensitive \\
ARF & 9,154 & 7.8\% & $+$0.036 & $6.5 \times 10^{-4}$ & *** & Sensitive \\
bHLH & 7,925 & 7.4\% & $+$0.043 & $1.5 \times 10^{-4}$ & *** & Sensitive \\
SBP & 8,501 & 3.7\% & $+$0.063 & $6.9 \times 10^{-9}$ & *** & Sensitive \\
\bottomrule
\end{tabular}
\end{table}

\begin{table}[!ht]
\centering
\caption{\textbf{Family-specific disruption predicts allelic expression magnitude.} For each TF family, median $|\log_2(\text{B73/NAM})|$ is compared between genes where that family's motifs are disrupted vs.\ not disrupted. Only families with $p < 0.05$ (Mann-Whitney) are shown.}
\label{tab:family_disruption}
\small
\begin{tabular}{lrrrrl}
\toprule
\textbf{Family} & \textbf{With} & \textbf{Without} & \textbf{Diff} & \textbf{$p$} & \textbf{Sig.} \\
\midrule
ERF & 0.555 & 0.481 & +0.074 & $6.0 \times 10^{-22}$ & *** \\
LBD & 0.565 & 0.509 & +0.056 & $8.2 \times 10^{-14}$ & *** \\
bZIP & 0.561 & 0.505 & +0.056 & $1.7 \times 10^{-11}$ & *** \\
bHLH & 0.559 & 0.507 & +0.052 & $1.5 \times 10^{-11}$ & *** \\
TCP & 0.556 & 0.512 & +0.044 & $1.5 \times 10^{-6}$ & *** \\
SBP & 0.551 & 0.509 & +0.042 & $2.4 \times 10^{-8}$ & *** \\
ARF & 0.548 & 0.511 & +0.037 & $9.7 \times 10^{-7}$ & *** \\
NAC & 0.543 & 0.506 & +0.037 & $3.1 \times 10^{-6}$ & *** \\
C2H2 & 0.536 & 0.508 & +0.028 & $3.0 \times 10^{-3}$ & ** \\
GATA & 0.547 & 0.522 & +0.025 & $4.3 \times 10^{-3}$ & ** \\
WRKY & 0.536 & 0.521 & +0.015 & $1.4 \times 10^{-2}$ & * \\
\bottomrule
\end{tabular}
\end{table}

% ============================================================
% FIGURES
% ============================================================

\begin{figure}[H]
    \centering
    \includegraphics[width=\textwidth]{figures/fig1_redundancy_overview.png}
    \caption{\textbf{Motif redundancy at bQTL positions.} (A)~Distribution of redundancy counts: number of additional same-family TF binding motifs within $\pm$200~bp of each bQTL. Red dashed line marks sole motifs (4.2\%). (B)~$|\log_2(\text{B73/NAM})|$ allelic expression by redundancy level. Box plots show median and IQR. (C)~Per-family Spearman correlation between redundancy and ASE. Dark red = $p < 0.001$; medium red = $p < 0.01$; light red = $p < 0.05$; gray = not significant.}
    \label{fig:redundancy}
\end{figure}

\begin{figure}[H]
    \centering
    \includegraphics[width=\textwidth]{figures/fig2_combinatorial_two_tier.png}
    \caption{\textbf{Combinatorial disruption and two-tier regulatory architecture.} (A)~Number of TF families with motifs disrupted at promoter bQTL vs.\ allelic expression magnitude ($|\log_2(\text{B73/NAM})|$). Spearman $\rho = +0.063$ ($p = 1.2 \times 10^{-13}$). (B)~Two-tier architecture: families sorted by their redundancy--ASE correlation. Green = buffered (negative $\rho$, more motifs reduce ASE); red = sensitive (positive $\rho$, disruption amplifies ASE). Dark shading = statistically significant ($p < 0.05$).}
    \label{fig:combinatorial}
\end{figure}

\begin{figure}[H]
    \centering
    \includegraphics[width=\textwidth]{figures/fig3_disruption_fraction.png}
    \caption{\textbf{Disruption fraction vs.\ allelic expression for buffered and sensitive families.} (A)~Buffered families (MYB, MYB\_related, Dof): higher disruption fraction (fewer alternative motifs) correlates with larger ASE. Binned medians (large markers) show the dose-response relationship. (B)~Sensitive families (WRKY, bHLH, SBP, bZIP): the relationship is reversed or absent, reflecting non-additive regulatory logic.}
    \label{fig:disruption}
\end{figure}

% ============================================================
\bibliographystyle{unsrtnat}
\begin{thebibliography}{15}

\bibitem[Engelhorn et al.(2025)]{engelhorn2025}
Engelhorn, J., Savadel, S.D., Banf, M., et al.
\newblock A pan-cistrome of maize identifies allelic cis-regulatory mechanisms of phenotypic divergence.
\newblock \textit{Nature Genetics}, 57:2313--2322, 2025.

\bibitem[de Boer et al.(2020)]{deboer2020}
de Boer, C.G., Vaishnav, E.D., Sadeh, R., et al.
\newblock Deciphering eukaryotic gene-regulatory logic with 100 million random promoters.
\newblock \textit{Nature Biotechnology}, 38:56--65, 2020.

\bibitem[Jores et al.(2021)]{jores2021}
Jores, T., Tonnies, J., Wrightsman, T., et al.
\newblock Synthetic promoter designs enabled by a comprehensive analysis of plant core promoters.
\newblock \textit{Nature Plants}, 7:842--855, 2021.

\bibitem[Gotea et al.(2010)]{gotea2010}
Gotea, V., Visel, A., Westlund, J.M., et al.
\newblock Homotypic clusters of transcription factor binding sites are a key component of human promoters and enhancers.
\newblock \textit{Genome Research}, 20(5):565--577, 2010.

\bibitem[Spitz and Furlong(2012)]{spitz2012}
Spitz, F. and Furlong, E.E.M.
\newblock Transcription factors: from enhancer binding to developmental control.
\newblock \textit{Nature Reviews Genetics}, 13:613--626, 2012.

\bibitem[Marand et al.(2021)]{marand2021}
Marand, A.P., Chen, Z., Gallavotti, A., Schmitz, R.J.
\newblock A cis-regulatory atlas in maize at single-cell resolution.
\newblock \textit{Cell}, 184(11):3041--3055, 2021.

\bibitem[Hwang et al.(2014)]{hwang2014}
Hwang, S.M., Kim, D.W., Woo, M.S., et al.
\newblock Functional characterization of Arabidopsis HsfA6a as a heat-shock transcription factor under high salinity and dehydration conditions.
\newblock \textit{Plant, Cell \& Environment}, 37(5):1202--1222, 2014.

\bibitem[Ye et al.(2018)]{ye2018}
Ye, Y., Wu, K., Chen, J., et al.
\newblock OsSND2, a NAC family transcription factor, is involved in secondary cell wall biosynthesis through regulating MYBs expression in rice.
\newblock \textit{Rice}, 11:36, 2018.

\bibitem[Sun et al.(2011)]{sun2011}
Sun, X., Feng, P., Xu, X., et al.
\newblock A chloroplast envelope-bound PHD transcription factor mediates chloroplast signals to the nucleus.
\newblock \textit{Nature Communications}, 2:477, 2011.

\bibitem[Ramegowda et al.(2017)]{ramegowda2017}
Ramegowda, V., Gill, U.S., Sivalingam, P.N., et al.
\newblock GBF3 transcription factor imparts drought tolerance in Arabidopsis thaliana.
\newblock \textit{Scientific Reports}, 7:9148, 2017.

\bibitem[Savadel et al.(2021)]{savadel2021}
Savadel, S.D., Hartwig, T., Turpin, Z.M., et al.
\newblock The native cistrome and sequence motif families of the maize ear.
\newblock \textit{PLoS Genetics}, 17(8):e1009689, 2021.

\bibitem[PlantTFDB(2017)]{planttfdb}
Jin, J., Tian, F., Yang, D.C., et al.
\newblock PlantTFDB 4.0: toward a central hub for transcription factors and regulatory interactions in plants.
\newblock \textit{Nucleic Acids Research}, 45(D1):D1040--D1045, 2017.

\end{thebibliography}

\end{document}
